% Abstract Page
\vspace{1in}
\chapter*{Abstract}


IT procurement is an important activity in organizations involved in the purchase, licensing, renewal and risk assessment of IT resources. The conventional procurement process involves manual checks, spreadsheets, emails and multiple systems leading to slow down, procurement of the same items, unused licences and blind spots about vendor risk.A variety of techniques including rule-based, document intensive, OCR, machine learning, and approval automation techniques can be employed to improve procurement processes but individual techniques typically focus on one or more stages in the procurement cycle but are limited in providing an end-to-end view of procurement activities. As a result, there is a need to develop a system that helps procure IT resources by orchestrating multiple procurement activities with minimal human intervention.
The application of Agentic AI is appropriate in this scenario as multiple specialised agents can be deployed to execute and coordinate various functions
In this paper, an Agentic AI-Based IT Procurement Intelligence Platform is designed based on a modular event driven monolithic architecture. The architecture comprises document and vendor intelligence agents, approval and inventory agents, contract and risk intelligence agents and the platform core services which operate on the principles of Apache Kafka, FastAPI, PostgreSQL, Redis and React technologies. Document processing techniques including OCR, classification and information extraction are deployed for purchase order, invoice and vendor quotation processing while machine learning techniques are applied in detecting duplicate purchases, analysing utilisation of licences and assessing vendor risk. The system integrates information and automates a variety of procurement activities with human agent involvement for low confidence and high priority items. The expected outcome of the developed system is to reduce the manual processing of procurement activities, improve procurement process efficiency, reduce unwarranted procurement expenditures and provide insights on vendor and procurement risks.